\section{What the method isolates}\label{sec:discussion}
Prime signs are independent, but their products satisfy exact identities
at arbitrarily large distances. The proof uses both facts. It counts the
homogeneous relation spaces created by multiplication, then conditions on
the prime environment that makes a dependency graph available. The
three-term profile \eqref{eq:master-profile}, and its full-value extension
\eqref{eq:absolute-profile}, are the arithmetic outputs. The independent
Poisson field is obtained only after the separate marginal, local-overlap
and prime-support estimates have been verified.

\subsection{A class of observables, rather than one count}
The dictionary theorem separates local word combinatorics from nonlocal
arithmetic. Changing the prescribed signs does not change a valuation
matrix. Prescribing their absolute orientation removes at most two parity
conditions in a pair of windows; the host count controls the additional
relations. Consequently the same profile handles a growing dictionary and
retains the identity and position of every rare word. Its overlap condition
holds for most dictionaries in the stated regime, and explicit marker
families make it zero. Short suffix--prefix matches can instead create
local clusters independently of any multiplicative relation; such
dictionaries require a separate analysis.

Long runs provide two complementary examples. Counting left-maximal starts
gives a Poisson field with geometric excesses and fair signs. Counting all
constant windows retains the clusters, and the field contracts to the
explicit compound-Poisson law of \cref{cor:window-clusters}. Both statements
come from the same comparison, not from two unrelated moment arguments.
The asymmetric almost-sure bounds in \cref{cor:as-longest} use its summable
quantitative errors; they do not require independence across scales.

The boundary is a different arithmetic regime. Its initial positive
cylinder imposes $\pi(L)$ independent prime constraints, whereas other
microscopic starts have an additional rank cost. The competition between
$2^{-\pi(L)}$ and $M2^{-L}$ therefore produces a change in the location
and sign of a rare long run. The two-clock theorem retains more than the
limiting source weights: it gives a prelimit lattice law, with prime-rank
overshoot at the boundary and integer overshoot in the bulk. The logistic
formula is a consequence of normalizing these two masses. The substantive
inputs are the microscopic rank surplus, the relative marked bulk kernel
and the negligible region between them.

\subsection{Arithmetic and probabilistic limits of the estimates}
Three terms reach $N^{5/3}$ at criticality in the present upper bound.
Improving only the terminal kernel energy would therefore not improve the
total exponent. The rational families in
\eqref{eq:rational-lower-dyadic}--\eqref{eq:rational-lower-macro} provide
lower exponents $4/3$ and $3/2$. Sharpening this gap is an arithmetic
question. It is distinct from improving the leading subexponential error,
which is governed by exceptional supports and shared prime coordinates.
For growing dictionaries the relevant quantity is not the full relation
sum, but its capped version. Mutual exclusivity at one site gives a
probability ceiling, and \cref{prop:capped-profile} preserves it through
the terminal count. This removes the terminal restriction on dictionary
size without improving the uncapped $5/3$ exponent. The remaining term
$NQ_B^{2/3}$ gives the sufficient exponent $1/2$ in
\eqref{eq:dictionary-critical}. A further improvement would need to
address this nonterminal contribution or use more information about the
dictionary; it does not follow from another terminal refinement alone.

A similar economy occurs in the Diophantine step. All component variables
already lie in bounded polynomial height. Localizing their squarefree
parts to the primes in the coefficient and the root differences reduces
the required count to two factors and a Pell equation. A height-free bound
for all integral points would be stronger, but is not needed. The proof
uses the hypotheses its application actually supplies.

The scalar, aggregated and spatially labelled comparisons retain different
information and therefore incur different costs; the table in
\cref{sec:conditional-laws} records their sufficient budgets. The local
resolution condition is separate: the absolute error must be negligible
relative to the probability of the event being resolved.

\subsection{Further questions}
A natural extension is a dictionary with substantial short overlaps.
The arithmetic profile still bounds distant interactions, but a new
cluster analysis is needed to transfer the local clump structure to
this multiplicative model. Such compound-Poisson descriptions for
Markov word families are developed in
\cite{ReinertSchbath1998,RoquainSchbath2007}; the problem here is their
compatibility with the nonlocal arithmetic dependence.
The constant-window compound-Poisson law gives a model case; it is not a
classification of periodic or interacting patterns. A second question is
the full one-factor growing-intensity range for the spatially labelled
field. This would require a sharper treatment of dependencies between
several sites, beyond the process comparison used here. Neither question
is an assumption of any theorem in this article.

Finally, the run lengths and their phase-indexed extremes remain on an
integer lattice. The separate process article
\cite{PoulyLatticePoissonFlows2026} develops records, occupation laws and
scale flows from the signed field. Its intrinsic integer observables are
distinct from its continuous-time embeddings and auxiliary phase removal,
which use additional randomization. None supplies an exact continuous
finite-$N$ Markov property for the multiplicative sequence.
